\section{实验背景与数据集}

\subsection{实验背景}
本实验构建了一个三层全连接神经网络。
输入处理：将 MNIST 数据集中的 28×28 灰度图像展平为 784 维向量，并在输入层添加偏置项（全1列），得到 785 维输入。

网络结构：第一层全连接层将输入从 785 维映射到 256 维，后接 ReLU 激活函数；第二层全连接层将 256 维映射到 128 维，后接 ReLU 激活函数；第三层全连接层将 128 维映射到 10 维，后接 Softmax 函数输出各类别概率。

权重初始化：采用 He 初始化方法，各层权重按 $\mathrm{scale} = \sqrt{2.0 / \mathrm{input\_dim}}$ 设置随机初始值。

模块实现：将矩阵乘法、ReLU 激活、Softmax 输出、交叉熵损失分别封装为独立类，每个类均实现 forward 前向传播和 backward 反向传播方法，便于梯度计算与参数更新。
\subsection{数据集}

本实验使用 \textbf{MNIST 手写数字数据集}：

\begin{itemize}
\item 训练集：60000 张图片
\item 测试集：10000 张图片
\item 图片大小：28×28（灰度图）
\item 分类类别：10 类（数字 0–9）
\end{itemize}

\subsection{数据预处理}

\subsubsection{归一化处理}

\[
x = x / 255.0
\]

\subsubsection{标签处理}

标签转为 One-hot 编码。

\subsubsection{输入展平}

\[
28 \times 28 \rightarrow 784
\]

\subsubsection{添加偏置项}

输入层添加偏置项（+1）。

